\documentclass[12pt,oncampus]{article}
\usepackage{stfx_syllabus}

\begin{document}

\stfxsyllabusheader{ST: Digital Methods for Historical Documents}{RELS 398.10 / ENGL 398.10}{Fall 2026}

\section*{Welcome to RELS 398 / ENGL 398!}

\instructorinfo{Dr.\ Ken Penner}{Dr.\ Penner / Professor Penner}{he/him/his}{kpenner@stfx.ca}{Nicholson Tower 503 \newline Book via Outlook Bookings link on Kwe' (Moodle)}

\textbf{Class Meetings:} Tuesdays 8:30 AM, Wednesdays 10:30 AM, Fridays 9:30 AM (Timeblock K1/K2/K3) in Mulroney Hall 3032 (MULH 3032)

\textbf{Contacting your instructor:} If you have any questions, you can contact me by email. Please include the course code (``RELS 398'' or ``ENGL 398'') and the topic in the subject line and sign the message with your full name and student number. I will try to respond to your email within 24 hours (excluding weekends, holidays, and exceptional circumstances).

\section*{StFX University Land Acknowledgement}

StFX is located in Mi’kma’ki, the unceded ancestral territory of the Mi'kmaw people. The Mi'kmaw name for Antigonish is Nalikitquniejk, meaning ‘place where branches are torn off.’

Ms't wiaqpulti'kl ankukamkewe'l. We are all treaty people.

\section*{Course Overview}

A practical, hands-on course on digital methods for processing historical texts and ancient manuscripts. Students develop practical skills in cataloguing, transcription, textual criticism, linguistic markup, translation, annotation, and digital publication. The course bridges traditional humanities scholarship with modern computational methods, preparing students for digital humanities research across disciplines.

The course follows a cumulative \textbf{learning-by-project} design: in Week~1, each student selects a primary manuscript witness, and each weekly assignment throughout the term develops one layer of a complete, TEI-conformant digital scholarly edition of that passage---spanning cataloguing and metadata, diplomatic transcription, double-keying reconciliation, orthographic normalization, rendering for web and print, morphological tagging, named-entity markup, translation alignment, and scholarly commentary, culminating in a print-ready digital edition portfolio.

To foster rigorous critical judgment, standard course weeks follow a \textbf{three-step comparative cycle}:
\begin{enumerate}[leftmargin=*,noitemsep,topsep=2pt]
  \item \textbf{Tuesday (Manual):} Transcribing, regularizing, or parsing by hand in a plain text editor to produce an authoritative scholarly \textbf{ground truth}.
  \item \textbf{Wednesday (Deterministic):} Applying specialized, rule-based software tools (e.g., Transkribus HTR models, XSLT stylesheets, CLTK pipelines).
  \item \textbf{Friday (AI-Assisted):} Evaluating modern LLMs and AI coding assistants against your ground truth to measure character and word error rates (CER/WER), audit hallucinations, and adjudicate editorial decisions.
\end{enumerate}

\subsection*{Course Outcomes}

By the end of this course, you should be able to:
\begin{itemize}
  \item Understand and apply cataloguing techniques and metadata standards (RDF, TEI header) for historical manuscripts.
  \item Accurately transcribe text from manuscript images using manual, semi-automated, and automated OCR/HTR platforms (e.g., Transkribus, THREAD).
  \item Render a TEI edition for readers on the web and in print, by hand, with XSLT, and with AI assistance.
  \item Apply TEI XML schema, linguistic markup, morphological tagging, and named-entity markup to historical documents.
  \item Produce interlinear glosses, translation alignments, scholarly commentary, and print-ready digital editions.
\end{itemize}

\subsection*{Course Materials \& Tools}

All reading assignments, software toolkits, and text/image repositories are provided electronically via Kwe' (Moodle) or open-access digital repositories (TEI Guidelines, GitHub, Transkribus, CETEIcean, CLTK).

\subsubsection*{Technology Requirements}
\begin{itemize}
  \item An internet-enabled laptop computer (Windows, macOS, or Linux) to bring to lab sessions.
  \item Git, Python 3, Jupyter Notebooks, and Visual Studio Code (VS Code) with XML extension (guidance provided in Week 1).
  \item Accounts on GitHub, Transkribus, and THREAD (setup in class).
\end{itemize}

\section*{Course Evaluation}

The following table details the assessment components for RELS 398 / ENGL 398.

\begin{stfxevaluation}

\evalrow{Cataloguing \& Metadata}{
  \begin{evalitemlist}
    \item Manuscript description, RDF, and TEI header metadata assignment.
  \end{evalitemlist}
}{3.5}

\evalrow{Transcription}{
  \begin{evalitemlist}
    \item Manual and automated manuscript image transcription (Transkribus \& THREAD).
  \end{evalitemlist}
}{10.5}

\evalrow{Normalization \& Regularization}{
  \begin{evalitemlist}
    \item Orthographic regularization and AI normalization assignment.
  \end{evalitemlist}
}{3.5}

\evalrow{Linguistic Markup}{
  \begin{evalitemlist}
    \item TEI XML morphological tagging and POS tagging assignment.
  \end{evalitemlist}
}{3.5}

\evalrow{Named-Entity Markup}{
  \begin{evalitemlist}
    \item People, places, organizations and dates tagged in TEI and linked to authority records.
  \end{evalitemlist}
}{3.5}

\evalrow{Interlinear Glossing}{
  \begin{evalitemlist}
    \item Word-level alignment and interlinear gloss assignment.
  \end{evalitemlist}
}{3.5}

\evalrow{Translation Alignment}{
  \begin{evalitemlist}
    \item Parallel text alignment and translation commentary.
  \end{evalitemlist}
}{14}

\evalrow{Scholarly Commentary}{
  \begin{evalitemlist}
    \item Detailed commentary on assigned manuscript passage features.
  \end{evalitemlist}
}{14}

\evalrow{Edition Introduction}{
  \begin{evalitemlist}
    \item Introductory essay on manuscript tradition and editorial principles.
  \end{evalitemlist}
}{7}

\evalrow{Print-Ready Publication}{
  \begin{evalitemlist}
    \item Final formatted digital edition publication package, building on the Week 7 web and print rendering.
  \end{evalitemlist}
}{3.5}

\evalrow{Comprehensive Digital Edition Portfolio}{
  \begin{evalitemlist}
    \item Assembled TEI-conformant edition submitted during the December exam period.
  \end{evalitemlist}
}{3.5}


\evalrow{Interactive Lessons \& Quizzes}{
  \begin{evalitemlist}
    \item Lesson and weekly quizzes completed before class (flipped classroom). Quizzes not attempted count as zero.
  \end{evalitemlist}
}{30}

\evaltotalrow{100}

\end{stfxevaluation}

\section*{Important Information for Students}

\subsection*{Academic Integrity \& AI Policy}

Academic honesty is essential in software code, XML markup, and scholarly analysis. Modern AI coding assistants (such as GitHub Copilot, ChatGPT, Claude) are introduced in this course as topic subjects and tools. However, all submitted XML markup, annotations, and commentaries must represent your own verified work. Any AI tool usage must be documented in your project repository commit log.

\subsection*{Class Attendance \& Special Dates}

Laboratory work requires active hands-on participation during scheduled class blocks. Note that on \textbf{Friday, September 25, 2026}, class will be conducted via \textbf{remote delivery} due to conference travel. On \textbf{Friday, December 4, 2026}, class follows a \textbf{Wednesday timetable} and will meet at \textbf{10:30 AM} in MULH 3032.

\end{document}
